%%%%%%%%%%%%%%%%%%%%%%%%%%%%%%%%%%%%%%%%%%%%%%%%%%%%%%%%%%%%%%%%%%%%%%%%%%%
% Lehreinheit 8: Quantitative Methoden und Statistik
% Quelle: WuZ-Buch, Kap. 15 (LE 8), Kap. 12d (Becher-Monty-Hall,
%         Beats 2--6), Kap. 12f (Feldstudie, nur Zahlen aus dem Buch),
%         Kap. 5b (Satz von Bayes, Horwich), Abschnitt 10d (p-Hacking)
% Staunmoment: Monty Hall mit drei Bechern (optional Chop Cup)
% Kernbegriff: Bedingte Wahrscheinlichkeit
%%%%%%%%%%%%%%%%%%%%%%%%%%%%%%%%%%%%%%%%%%%%%%%%%%%%%%%%%%%%%%%%%%%%%%%%%%%
\nsl{Bis hierher hat das Semester gefragt, wie man eine Frage stellt, wie
man Quellen prüft und wie man ein Design baut, das sich selbst widerlegen
kann. Diese Einheit führt das erste große Werkzeug ein, mit dem man Daten
befragt: die Statistik. Sie beginnt mit drei Bechern und einer
Abstimmung, bei der die Mehrheit fast immer danebenliegt. Das ist der
eigentliche Grund, warum es Statistik gibt: Unser Gefühl für
Wahrscheinlichkeiten ist schlecht geeicht, und zwar nicht zufällig,
sondern systematisch. Die Einheit führt von dort zu den Grundbegriffen der
beschreibenden und der schließenden Statistik, zum p-Wert und zu dem, was
er \emph{nicht} sagt, und zur bayesianischen Frage, wie viel ein Befund
eigentlich beweist.
\newpage}
\os{\begin{description}
\item[{\rot\bf Heute}]~\\[-\lblskip]
\bi
  \ite drei Becher, eine Kugel: wechseln oder bleiben?
  \ite bedingte Wahrscheinlichkeit: warum das Sichtbare täuscht
  \ite Daten beschreiben: Skalen, Lage, Streuung, Ausreißer
  \ite Daten prüfen: Signifikanztest, p-Wert, Effektgröße, Power
  \ite Bayes: was ein Befund beweist und was nicht
\ei
\end{description}
\newpage}

\begin{vorbereitung}[Monty Hall mit drei Bechern (optional: Chop Cup)]
\textbf{Eine Woche vorher}
\begin{checkliste}
  \item Für die Vorführung: drei \emph{undurchsichtige}, gleich aussehende
    Becher und eine gut sichtbare Kugel (Schaumstoffball, Walnuss). Die
    Becher müssen von hinten im Saal zu sehen sein; bei großen Räumen eine
    Dokumentenkamera einplanen.
  \item Für das Nachspielen in Paaren: je Paar drei undurchsichtige
    Pappbecher, eine Kugel (Papierkugel, Erbse, Münze) und ein Würfel.
    Bei 30 Studierenden 15 Sets plus drei Reserve. Pappbecher und Würfel
    gibt es im Haushalts- bzw.\ Spielwarenhandel.
  \item Optional als Staunmoment des Vortragenden: eine
    \emph{Chop-Cup-Routine} (ca.\ 2 Min.). Das Requisit gibt es im
    Zauberfachhandel. Methode und Einstudierung nach der Anleitung des
    Herstellers bzw.\ der Standardliteratur; hier bewusst nicht
    beschrieben. Die Routine muss sitzen, bevor sie vor Publikum gezeigt
    wird: mindestens zehn Durchläufe vor dem Spiegel und einer vor
    einer Testperson.
  \item Die drei Fälle und die Gegenprobe mit dem blinden Moderator
    (Buch, Kap.~12d, Beats~2--6) so oft mit echten Bechern durchspielen,
    dass sie ohne Folie frei erklärt werden können.
  \item Für den p-Hacking-Kurztest (Münzwurf-\enquote{Hellseher}) die
    Studierenden per Mail bitten, eine Münze mitzubringen; 20 Ersatzmünzen
    einstecken.
\end{checkliste}

\textbf{Handout und Übungsblatt}
\begin{checkliste}
  \item Handout \enquote{Quantitative Methoden und Statistik} sowie
    Übungsblatt LE~8 Teil~A und Teil~B je einmal pro Person kopieren
    (Druckvorlage \texttt{handouts.pdf}, Seiten zu LE~8). Teil~A liegt zu
    Beginn aus und enthält die Strichliste für das Monty-Hall-Spiel und
    den vollständigen fiktiven Datensatz ($n=50$) für die
    Datensatz-Analyse. Teil~B (Zeitungsmeldung, Hausaufgaben) fragt nach
    der Zahl der Vergleiche und nimmt damit die
    Münzwurf-\enquote{Hellseher}-Demo vorweg; es wird erst in der
    Transferphase ausgeteilt.
  \item Je Kleingruppe ein Blatt Millimeterpapier (A4) für Histogramm und
    Boxplot; Taschenrechner erlauben, Tabellenkalkulation optional.
  \item Für die Transferphase zwei bis drei aktuelle Zeitungsmeldungen
    mitbringen, die eine Studie mit \enquote{signifikant} oder
    \enquote{bewiesen} zitieren. Gibt es gerade keine passende Meldung,
    reicht die fiktive Beispielmeldung auf dem Übungsblatt (Teil~B,
    Aufgabe~3).
\end{checkliste}

\textbf{Am Tag selbst}
\begin{checkliste}
  \item Tafel in drei Felder teilen: \emph{Abstimmung} (wechseln /
    bleiben / egal), \emph{Kursdaten} (Summen der Strichlisten),
    \emph{Fragen}.
  \item Becher-Sets vor Beginn auf die Tische verteilen, aber gestapelt
    lassen, damit niemand vorher experimentiert.
  \item Falls die Chop-Cup-Routine gezeigt wird: das präparierte Requisit
    danach sichtbar wegräumen und das Monty-Hall-Spiel mit gewöhnlichen
    Bechern spielen, die man den Studierenden in die Hand gibt. Ein Spiel,
    dessen Requisiten man misstraut, liefert keine brauchbaren Daten.
\end{checkliste}

\textbf{Risiken und Plan B.} Ein Teil des Kurses kennt das
Ziegenproblem. Das schadet nicht, es ist ein Datum: Vor der
Abstimmung fragen, wer das Rätsel schon kennt, und beide Gruppen
getrennt auszählen. Die Feldstudie des Buchs zeigt, wie stark Vorwissen
wirkt (Kap.~12f; siehe Abschnitt \enquote{Effektgröße und Power}).
Laufen die Paarspiele schleppend, nach 15 Durchgängen abbrechen; der
Kurs-Pool reicht dann immer noch für die Rechnung. Wird die
Chop-Cup-Routine nicht gezeigt, beginnt die Einheit direkt mit dem
Monty-Hall-Spiel; der Staunmoment entsteht dort aus dem Abstand zwischen
Abstimmung und Strichliste.
\end{vorbereitung}

\begin{ablauf}
\begin{description}
\item[0--10 · Staunmoment.] Optional die Chop-Cup-Routine. Dann
  Monty Hall live mit drei Bechern; Abstimmung wechseln / bleiben / egal.
  Ergebnis an die Tafel, ohne Kommentar.
\item[10--20 · Nachspielen.] 20--30 Durchgänge in Paaren nach festen
  Regeln (Würfel legt die Kugel, Strategie wechselt jede Runde),
  Strichliste auf dem Übungsblatt Teil~A (Aufgabe~1). Summen an die
  Tafel.
\item[20--50 · Theorie-Input.] Auflösung über die drei Fälle und den
  blinden Moderator; deskriptiv vs.\ inferenz; Skalenniveaus; Lage und
  Streuung; Normalverteilung; Signifikanztest, p-Wert und was er
  \emph{nicht} sagt; Effektgröße und Power; Bayes und der Beweiswert
  riskanter Vorhersagen.
\item[50--60 · Warnung.] Münzwurf-\enquote{Hellseher} als
  p-Hacking-Kurzdemo; Ausblick auf den Vorhersage-Trick in LE~10.
\item[60--80 · Datensatz-Analyse.] Kleingruppen beschreiben und
  visualisieren den fiktiven Datensatz (Übungsblatt Teil~A, Aufgabe~2).
\item[80--90 · Transfer.] Übungsblatt Teil~B austeilen. Eine
  Zeitungsmeldung gemeinsam methodisch lesen (Aufgabe~3); Auftrag zur
  Kurzeinschätzung (Aufgabe~4).
\end{description}
\end{ablauf}

%==============================================================================
\section{Staunmoment: Drei Becher}
%==============================================================================
\begin{staunbox}{Wechseln oder bleiben?}
Unter einem von drei Bechern liegt eine Kugel. Sie zeigen auf einen.
Ich weiß, wo die Kugel ist, und hebe einen \emph{anderen}, leeren Becher
hoch.

\emph{Bleiben Sie bei Ihrem Becher -- oder wechseln Sie zum anderen?
Oder ist es egal?}
\end{staunbox}
\os{\newpage}

\begin{description}
\item[{\rot\bf Abstimmung}]~\\[-\lblskip]
\bi
  \ite Hand hoch: \textbf{bleiben} -- \textbf{wechseln} -- \textbf{egal}
  \ite Wer kennt das Rätsel schon? (getrennt zählen)
  \ite Wie sicher sind Sie? Schätzen Sie in Prozent.
\ei
\nsl{Die Abstimmung wird an der Tafel festgehalten und vorerst nicht
kommentiert. In der Regel stimmt die Mehrheit für \enquote{egal}, oft mit
hoher Sicherheit. Das Rätsel ist als \emph{Ziegenproblem} bekannt
geworden; es geht auf \textcite{mhall_selvin1975} zurück und wurde durch
die Kolumne von \textcite{mhall_vossavant1990} berühmt, auf die zahlreiche
Leserbriefe mit der Behauptung antworteten, die Lösung sei falsch.
\textcite{mhall_krauss2003} zeigen, wie hartnäckig die Fehlintuition ist
und welche Darstellungen helfen.}
\os{\newpage}

\item[{\rot\bf Nachspielen: die Regeln}]~\\[-\lblskip]
\bi
  \ite in Paaren: eine Person ist \emph{Spielleitung}, eine spielt
  \ite Spielleitung würfelt verdeckt: 1--2 Becher A, 3--4 B, 5--6 C
  \ite Spielleitung kennt die Lage und öffnet immer einen \emph{leeren}
    Becher, nie den gewählten
  \ite ungerade Runden: \textbf{bleiben}, gerade Runden: \textbf{wechseln}
  \ite 20--30 Runden, Strichliste auf dem Übungsblatt Teil~A
\ei
\nsl{Die Regeln sind bereits ein kleines Forschungsdesign. Der Würfel
sorgt für eine zufällige Lage der Kugel, damit die Spielleitung nicht
unbewusst Lieblingsbecher bevorzugt. Die feste Abfolge von Bleiben und
Wechseln verhindert, dass die spielende Person ihre Strategie nach Gefühl
wählt und damit die Vergleichbarkeit der beiden Strategien zerstört. Hat
die spielende Person zufällig den Becher mit der Kugel gewählt, entscheidet
ein zweiter Wurf (gerade oder ungerade), welchen der beiden leeren Becher
die Spielleitung öffnet.}
\os{\newpage}

\item[{\rot\bf Was sagen die Daten?}]~\\[-\lblskip]
\bi
  \ite Kursdaten an der Tafel: Gewinne beim \textbf{Bleiben} und beim
    \textbf{Wechseln}
  \ite Passt das zu \enquote{egal}?
  \ite Was müsste herauskommen, wenn es egal wäre?
  \ite Wie groß darf die Abweichung sein, bevor wir \enquote{egal}
    aufgeben?
\ei
\nsl{Mit der letzten Frage ist die Einheit beim Kern der
schließenden Statistik angekommen, noch bevor das Wort fällt. Die Gruppe
hat eine Hypothese (\enquote{egal}, also Gewinnchance $\tfrac{1}{2}$ beim
Wechseln), eine Stichprobe (die Strichlisten) und eine Abweichung, deren
Gewicht sie beurteilen muss.}
\os{\newpage}
\end{description}

%==============================================================================
\section{Auflösung: bedingte Wahrscheinlichkeit}
%==============================================================================
\nsl{Die Auflösung folgt der Becher-Erklärung aus dem Buch (Kap.~12d,
Beats~2--6). Der Vorgang hat zwei Stufen: Zuerst liegt die Kugel zufällig
unter einem von drei Bechern, jeder mit Wahrscheinlichkeit $\tfrac{1}{3}$.
Dann öffnet die Spielleitung einen Becher, aber nicht irgendeinen: nie den
gewählten und nie den mit der Kugel. Die Fehlintuition \enquote{zwei
Becher, also fifty-fifty} beruht auf der stillen Annahme, die beiden
übrigen Becher seien gleichberechtigt. Das sind sie nicht. Der gewählte
Becher wurde nie angefasst; der andere ist der Überlebende einer Auswahl,
die mit Wissen getroffen wurde.}
\begin{description}
\item[{\rot\bf Die drei Fälle}]~\\[-\lblskip]
\begin{center}
\begin{tabular}{@{}llll@{}}
\toprule
\textbf{Kugel liegt \dots} & \textbf{Wahrsch.} & \textbf{Spielleitung öffnet} & \textbf{Wechseln} \\
\midrule
unter Ihrem Becher 1 & $\tfrac13$ & Becher 2 oder 3 & verliert \\
unter Becher 2 & $\tfrac13$ & zwingend Becher 3 & gewinnt \\
unter Becher 3 & $\tfrac13$ & zwingend Becher 2 & gewinnt \\
\bottomrule
\end{tabular}
\end{center}
\bi
  \ite in zwei von drei Fällen \emph{zwingt} die Lage der Kugel die
    Spielleitung, den richtigen Becher übrig zu lassen
  \ite Wechseln gewinnt mit $\tfrac23$, Bleiben mit $\tfrac13$
\ei
\os{\newpage}

\item[{\rot\bf Versteckte Information}]~\\[-\lblskip]
\bi
  \ite Ihr Becher hatte $\tfrac13$ und behält es: Die Spielleitung
    hätte ihn ohnehin nie geöffnet
  \ite die $\tfrac23$ der beiden anderen Becher sammeln sich auf dem
    einen, der übrig bleibt
  \ite die Spielleitung sagt kein Wort über die Kugel, aber die
    \emph{Art ihrer Auswahl} trägt Information
\ei
\os{\newpage}

\item[{\rot\bf Gegenprobe: der blinde Moderator}]~\\[-\lblskip]
\bi
  \ite die Spielleitung weiß \emph{nicht}, wo die Kugel ist, und öffnet
    zufällig einen anderen Becher
  \ite zufällig ist er leer: Das Bild ist exakt dasselbe wie vorher
  \ite jetzt, und nur jetzt, ist es tatsächlich egal
  \ite \textbf{dasselbe Sichtbare, zwei verschiedene Prozesse, zwei
    verschiedene Wahrscheinlichkeiten}
\ei
\nsl{Die Präzisierung aus dem Buch gehört dazu: Beim blinden Moderator
gilt fifty-fifty \emph{unter der Bedingung}, dass er zufällig einen leeren
Becher erwischt hat; deckt er die Kugel auf, ist die Runde hinfällig. Der
Unterschied zwischen beiden Spielvarianten liegt allein darin, ob der
zweite Schritt an die Lage der Kugel gekoppelt ist. Für die Statistik ist
das die wichtigste Lehre des Rätsels: Wer nur auf das Ergebnis schaut,
kann die beiden Fälle nicht unterscheiden. Wer fragt, wie die Daten
zustande gekommen sind, kann es.}
\os{\newpage}

\item[{\rot\bf Bedingte Wahrscheinlichkeit}]~\\
\defn{bedingtewkt}{Die \textbf{bedingte Wahrscheinlichkeit} $P(A\mid B)$
ist die Wahrscheinlichkeit von $A$ unter der Voraussetzung, dass $B$
eingetreten ist:
$P(A\mid B) = P(A\cap B)\,/\,P(B)$.}
\bi
  \ite Monty Hall: $P(\text{Kugel unter dem anderen Becher} \mid
    \text{ein leerer wurde geöffnet})$
  \ite wissende und blinde Spielleitung erzeugen \emph{verschiedene}
    Bedingungen, obwohl das Bild gleich ist
\ei
\os{\newpage}
\end{description}

\begin{merkbox}
Wahrscheinlichkeiten hängen nicht am Bild, sondern am Prozess, der das
Bild erzeugt hat. Vor jeder Auswertung steht deshalb die Frage: Wie sind
diese Daten zustande gekommen?
\end{merkbox}
\os{\newpage}

%==============================================================================
\section{Beschreiben und Schließen}
%==============================================================================
\nsl{Die Strichlisten der Paare sind ein kleiner Datensatz. Mit ihm lassen
sich zwei grundverschiedene Dinge tun. Man kann ihn \emph{beschreiben}:
Wie oft hat Wechseln in unserem Kurs gewonnen? Das ist deskriptive
Statistik; sie fasst vorliegende Daten zusammen und sagt nichts über
Fälle, die nicht erhoben wurden. Man kann aber auch über den Kurs hinaus
\emph{schließen}: Ist die Gewinnchance beim Wechseln in diesem Spiel
überhaupt, also in allen denkbaren Runden, größer als ein Halb? Das ist
Inferenzstatistik. Sie behandelt die erhobenen Daten als Stichprobe aus
einer Grundgesamtheit und muss deshalb mit Zufallsschwankungen rechnen.
Ein Standardlehrbuch für beide Teile ist \textcite{doeringbortz2016}.}
\begin{description}
\item[{\rot\bf Zwei Aufgaben der Statistik}]~\\[-\lblskip]
\begin{center}
\begin{tabular}{p{0.42\textwidth}p{0.42\textwidth}}
\toprule
\textbf{Deskriptive Statistik} & \textbf{Inferenzstatistik} \\
\midrule
beschreibt die vorliegenden Daten & schließt von der Stichprobe auf die Grundgesamtheit \\
Kennwerte, Tabellen, Grafiken & Tests, Konfidenzintervalle, Modelle \\
\enquote{In unserem Kurs gewann Wechseln in 64\,\% der Runden.} & \enquote{Wechseln gewinnt häufiger als in der Hälfte der Fälle.} \\
\bottomrule
\end{tabular}
\end{center}
\nsl{Die 64\,\% in der Tabelle sind ein Rechenbeispiel, kein Befund.}
\os{\newpage}

\item[{\rot\bf Skalenniveaus}]~\\[-\lblskip]
\begin{center}
\small
\begin{tabular}{@{}p{0.15\textwidth}p{0.30\textwidth}p{0.27\textwidth}p{0.16\textwidth}@{}}
\toprule
\textbf{Skala} & \textbf{sinnvoll ist \dots} & \textbf{Beispiel} & \textbf{Lagemaß} \\
\midrule
nominal & gleich / ungleich & Strategie (bleiben, wechseln) & Modus \\
ordinal & zusätzlich größer / kleiner & Staunen 1--5 & Median \\
intervall & zusätzlich Abstände & Temperatur in °C & Mittelwert \\
verhältnis & zusätzlich echter Nullpunkt & Ausschlag in mm & Mittelwert \\
\bottomrule
\end{tabular}
\end{center}
\bi
  \ite das Skalenniveau entscheidet, welche Rechnungen erlaubt sind
\ei
\nsl{Der häufigste Fehler betrifft Ratingskalen. Eine Antwort von~4 auf
einer Skala von~1 bis~5 ist nicht \enquote{doppelt so viel} wie eine~2,
und ob die Abstände zwischen den Stufen gleich groß sind, weiß niemand.
In der Praxis werden Mittelwerte aus Likert-Items trotzdem oft berichtet;
wer das tut, sollte es begründen und den Median mit angeben.}
\os{\newpage}
\end{description}

%==============================================================================
\section{Lage, Streuung und Ausreißer}
%==============================================================================
\nsl{Ein Datensatz lässt sich durch zwei Fragen grob erfassen: Wo liegt
seine Mitte, und wie weit streuen die Werte um sie herum? Für die Mitte
stehen Modus, Median und arithmetisches Mittel zur Verfügung. Der Median
ist der Wert, unter und über dem jeweils die Hälfte der Daten liegt; das
Mittel ist die Summe geteilt durch die Anzahl. Bei symmetrischen
Verteilungen liegen beide nahe beieinander. Bei schiefen Verteilungen oder
einzelnen Extremwerten wird das Mittel in Richtung der Ausreißer gezogen,
der Median bleibt stabil. Für die Streuung sind Spannweite,
Quartilsabstand und Standardabweichung gebräuchlich. Die
Standardabweichung misst die typische Entfernung der Werte vom Mittelwert
und ist ebenso ausreißerempfindlich wie dieser.}
\begin{description}
\item[{\rot\bf Lage: Wo liegt die Mitte?}]~\\[-\lblskip]
\bi
  \ite \textbf{Modus}: der häufigste Wert
  \ite \textbf{Median} $\tilde{x}$: die Hälfte liegt darunter, die Hälfte
    darüber
  \ite \textbf{Mittelwert} $\bar{x} = \frac{1}{n}\sum x_i$
  \ite Gedankenexperiment: zehn Gäste in einer Kneipe, ihr mittleres
    Einkommen -- dann kommt eine Milliardärin herein
\ei
\nsl{Durch eine einzige Person steigt der Mittelwert ins Absurde, der
Median ändert sich kaum. Welcher Kennwert \enquote{richtig} ist, hängt von
der Frage ab. Wer wissen will, was eine typische Person verdient, nimmt
den Median; wer wissen will, wie viel Geld insgesamt im Raum ist, braucht
den Mittelwert.}
\os{\newpage}

\item[{\rot\bf Streuung: Wie verschieden sind die Werte?}]~\\[-\lblskip]
\bi
  \ite \textbf{Spannweite}: Maximum minus Minimum
  \ite \textbf{Quartile} $Q_1$, $Q_3$: ein bzw.\ drei Viertel der Werte
    liegen darunter
  \ite \textbf{Quartilsabstand} $\mathit{IQR} = Q_3 - Q_1$: Breite der
    mittleren Hälfte
  \ite \textbf{Standardabweichung}
    $s = \sqrt{\tfrac{1}{n-1}\sum (x_i-\bar{x})^2}$
\ei
\os{\newpage}

\item[{\rot\bf Boxplot und Ausreißerregel}]~\\[-\lblskip]
\begin{center}
\resizebox{0.8\textwidth}{!}{%
\begin{tikzpicture}[font=\sffamily\small]
  \draw[->,rwugray] (0,0) -- (11,0) node[right]{Wert};
  \foreach \x in {0,2,...,10} \draw[rwugray] (\x,0.1) -- (\x,-0.1);
  \fill[rwucyan40!50] (3,0.8) rectangle (6,1.8);
  \draw[thick,rwuviolet] (3,0.8) rectangle (6,1.8);
  \draw[very thick,rwuviolet] (4.2,0.8) -- (4.2,1.8);
  \draw[thick,rwuviolet] (1,1.3) -- (3,1.3);
  \draw[thick,rwuviolet] (6,1.3) -- (8.2,1.3);
  \draw[thick,rwuviolet] (1,1.05) -- (1,1.55);
  \draw[thick,rwuviolet] (8.2,1.05) -- (8.2,1.55);
  \fill[rwuviolet] (10.2,1.3) circle (2.5pt);
  \node[above] at (3,1.8) {$Q_1$};
  \node[above] at (4.6,1.8) {Median};
  \node[above] at (6.2,1.8) {$Q_3$};
  \node[above] at (10.2,1.45) {Ausreißer};
  \draw[decorate,decoration={brace,mirror,amplitude=4pt},rwugray] (3,0.6) -- (6,0.6)
    node[midway,below=4pt]{$\mathit{IQR}$};
\end{tikzpicture}}
\end{center}
\bi
  \ite Box = mittlere Hälfte, Strich = Median, Antennen bis zum
    letzten unauffälligen Wert
  \ite Ausreißer nach \textcite{tukey1977eda}: mehr als
    $1{,}5\cdot\mathit{IQR}$ unter $Q_1$ oder über $Q_3$
\ei
\os{\newpage}

\item[{\rot\bf Was tun mit Ausreißern?}]~\\[-\lblskip]
\bi
  \ite zuerst prüfen: Messfehler, Tippfehler -- oder ein echter Fall?
  \ite \textbf{nie} still löschen, weil er \enquote{stört}
  \ite Regel \emph{vor} der Auswertung festlegen und berichten
  \ite im Zweifel beide Auswertungen zeigen: mit und ohne
\ei
\nsl{Ausreißer sind die erste Stelle, an der die Freiheit der Auswertung
beginnt. Wer erst nach dem Blick auf das Ergebnis entscheidet, ob ein
Extremwert drin bleibt, hat bereits einen Pfad im \enquote{Garten der
sich verzweigenden Wege} gewählt \parencite{gelmanloken2014forking}. Darauf
kommt LE~10 ausführlich zurück.}
\os{\newpage}

\item[{\rot\bf Daten zeigen, nicht verstecken}]~\\[-\lblskip]
\bi
  \ite Histogramm und Boxplot für metrische Daten, Balken für
    Häufigkeiten
  \ite Achsen beschriften, Einheit angeben, $n$ nennen
  \ite Fallstricke: abgeschnittene $y$-Achse, 3D-Effekte, Mittelwertbalken
    ohne Streuung
  \ite wenn möglich: die einzelnen Datenpunkte zeigen
\ei
\nsl{Eine abgeschnittene Achse ist nicht in jedem Fall unzulässig, sie
muss aber erkennbar sein. Ein Balkendiagramm, das nur zwei Mittelwerte
zeigt, verschweigt, wie stark sich die Gruppen überlappen. Je kleiner der
Datensatz, desto eher sollte man alle Punkte zeigen.}
\os{\newpage}
\end{description}

%==============================================================================
\section{Normalverteilung und Signifikanztest}
%==============================================================================
\nsl{Viele Messgrößen, die sich aus vielen kleinen, voneinander
unabhängigen Einflüssen zusammensetzen, verteilen sich annähernd
glockenförmig. Für diese Normalverteilung gilt eine einfache Faustregel:
Etwa 68\,\% der Werte liegen innerhalb einer Standardabweichung um den
Mittelwert, etwa 95\,\% innerhalb von zwei. Wichtiger als die
Verteilung einzelner Messwerte ist für die Inferenzstatistik eine andere
Tatsache: Mittelwerte und Anteile aus hinreichend großen Stichproben sind
näherungsweise normalverteilt, auch wenn die einzelnen Werte es nicht
sind. Darauf beruhen die meisten Standardtests.}
\begin{description}
\item[{\rot\bf Die Normalverteilung}]~\\[-\lblskip]
\begin{center}
\resizebox{0.75\textwidth}{!}{%
\begin{tikzpicture}[font=\sffamily\small,x=1.8cm,
  declare function={g(\x)=3*exp(-(\x)^2/2);}]
  \fill[rwucyan40!30,domain=-2:2,samples=60] (-2,0) -- plot (\x,{g(\x)}) -- (2,0) -- cycle;
  \fill[rwucyan40!70,domain=-1:1,samples=40] (-1,0) -- plot (\x,{g(\x)}) -- (1,0) -- cycle;
  \draw[very thick,rwuviolet,domain=-3.5:3.5,samples=100] plot (\x,{g(\x)});
  \draw[->,rwugray] (-3.8,0) -- (3.9,0);
  \foreach \x/\l in {-2/$\mu-2\sigma$,-1/$\mu-\sigma$,0/$\mu$,1/$\mu+\sigma$,2/$\mu+2\sigma$}
    {\draw[rwugray] (\x,0.08) -- (\x,-0.08); \node[below] at (\x,-0.08) {\l};}
  \node at (0,1.2) {ca.\ 68\,\%};
  \draw[<->,rwugray] (-2,3.3) -- node[above]{ca.\ 95\,\%} (2,3.3);
\end{tikzpicture}}
\end{center}
\os{\newpage}

\item[{\rot\bf Die Logik des Signifikanztests}]~\\[-\lblskip]
\bi
  \ite \textbf{Nullhypothese} $H_0$: Es gibt keinen Effekt
    (\enquote{Wechseln gewinnt mit $\tfrac12$})
  \ite Frage: Wie wahrscheinlich wären Daten, die \emph{mindestens so
    extrem} sind wie unsere, \emph{wenn $H_0$ stimmt}?
  \ite diese Wahrscheinlichkeit ist der \textbf{p-Wert}
  \ite ist $p$ kleiner als eine vorab gewählte Schwelle $\alpha$ (meist
    $0{,}05$), verwirft man $H_0$
\ei
\os{\newpage}

\item[{\rot\bf Rechenbeispiel: Kurs-Pool}]~\\[-\lblskip]
\bi
  \ite angenommen, der Kurs hat 150 Wechsel-Runden gespielt und 97 davon
    gewonnen (64,7\,\%)
  \ite unter $H_0$ (Chance $\tfrac12$) erwartet man 75
  \ite Binomialtest, zweiseitig: $p \approx 0{,}0004$
  \ite 95-\%-Konfidenzintervall für die Gewinnquote: etwa 57\,\% bis
    72\,\%
  \ite der Wert $\tfrac23$ liegt im Intervall, $\tfrac12$ nicht
\ei
\nsl{Die Zahlen sind ein konstruiertes Rechenbeispiel (exakter
Binomialtest gegen $\tfrac12$, Wilson-Intervall). In der Stunde wird mit
den echten Kursdaten gerechnet; eine Tabellenkalkulation oder ein
Online-Rechner für den Binomialtest genügt. Wichtig ist der letzte Punkt
der Folie: Das Konfidenzintervall sagt mehr als der p-Wert, weil es zeigt,
welche Gewinnquoten mit den Daten verträglich sind.}
\os{\newpage}

\item[{\rot\bf Zwei Arten, sich zu irren}]~\\[-\lblskip]
\begin{center}
\begin{tabular}{@{}lll@{}}
\toprule
 & \textbf{$H_0$ stimmt} & \textbf{$H_0$ stimmt nicht} \\
\midrule
\textbf{$H_0$ verworfen} & Fehler 1.~Art ($\alpha$) & richtig (Power $1-\beta$) \\
\textbf{$H_0$ beibehalten} & richtig & Fehler 2.~Art ($\beta$) \\
\bottomrule
\end{tabular}
\end{center}
\bi
  \ite $\alpha = 0{,}05$ heißt: In 5\,\% der Fälle ohne Effekt meldet der
    Test trotzdem einen
  \ite ein nicht signifikantes Ergebnis ist kein Beleg für \enquote{kein
    Effekt}
\ei
\os{\newpage}
\end{description}

%==============================================================================
\section{Was der p-Wert nicht sagt}
%==============================================================================
\nsl{Kaum ein statistischer Begriff wird so oft missverstanden wie der
p-Wert. Die American Statistical Association hat 2016 eine eigene
Stellungnahme veröffentlicht, um die gängigsten Fehldeutungen
richtigzustellen \parencite{wassersteinlazar2016asa}. Der Kern: Ein
p-Wert beschreibt, wie gut die Daten zu einem bestimmten Modell (der
Nullhypothese mit allen Zusatzannahmen) passen. Er sagt nichts darüber,
wie wahrscheinlich eine Hypothese ist, und nichts darüber, wie groß oder
wichtig ein Effekt ist.}
\begin{description}
\item[{\rot\bf Der p-Wert ist \emph{nicht} \dots}]~\\[-\lblskip]
\bi
  \ite die Wahrscheinlichkeit, dass $H_0$ wahr ist
  \ite die Wahrscheinlichkeit, dass das Ergebnis \enquote{Zufall} ist
  \ite ein Maß für die Größe oder Bedeutung eines Effekts
  \ite ein Beweis, weder für noch gegen eine Hypothese
  \ite unabhängig davon, wie viele Tests man gerechnet hat
\ei
\os{\newpage}

\item[{\rot\bf Warum \enquote{p < 0,05} kein Beweis ist}]~\\[-\lblskip]
\bi
  \ite $p$ fragt nach $P(\text{Daten} \mid H_0)$
  \ite wir wollen wissen: $P(H \mid \text{Daten})$
  \ite das ist wieder bedingte Wahrscheinlichkeit, und die Richtung
    der Bedingung lässt sich nicht umdrehen
  \ite was dazwischen fehlt, ist das Vorwissen über $H$ (der Prior)
\ei
\nsl{Die Verwechslung von $P(\text{Daten}\mid H_0)$ mit
$P(H_0\mid\text{Daten})$ ist ein naher Verwandter des Fehlers, den das
Ziegenproblem sichtbar macht: Man liest eine Bedingung falsch. Wie man
von der einen Wahrscheinlichkeit zur anderen kommt, regelt der Satz von
Bayes; dafür braucht man aber eine Vorab-Einschätzung der Hypothese. Ist
diese niedrig, bleibt die Hypothese auch nach einem kleinen p-Wert
unsicher. Das erklärt, warum überraschende Einzelbefunde so oft nicht
replizieren \parencite{ioannidis2005false,opensciencecollaboration2015}.}
\os{\newpage}
\end{description}

%==============================================================================
\section{Effektgröße und Power}
%==============================================================================
\nsl{Ob ein Effekt \enquote{signifikant} ist, hängt nicht nur von seiner
Größe ab, sondern auch von der Stichprobengröße. Mit sehr vielen Personen
wird jeder noch so kleine Unterschied signifikant, mit sehr wenigen
bleibt selbst ein großer unentdeckt. Deshalb gehört zu jedem Befund eine
\emph{Effektgröße}, die unabhängig von $n$ angibt, wie groß der
Unterschied ist. Für Mittelwertunterschiede ist Cohens $d$ üblich, der
Abstand der Mittelwerte in Einheiten der Standardabweichung
\parencite{cohen1988power}. Die \emph{Power} (Teststärke) ist die
Wahrscheinlichkeit, einen tatsächlich vorhandenen Effekt bestimmter Größe
auch zu finden. Sie wird \emph{vor} der Studie über eine Poweranalyse
festgelegt; daraus ergibt sich die nötige Stichprobengröße.}
\begin{description}
\item[{\rot\bf Effektgröße: Wie groß ist der Unterschied?}]~\\[-\lblskip]
\bi
  \ite Cohens $d = (\bar{x}_1 - \bar{x}_2)\,/\,s$: Abstand in
    Standardabweichungen
  \ite grobe Orientierung: $0{,}2$ klein, $0{,}5$ mittel, $0{,}8$ groß
    \parencite{cohen1988power}
  \ite immer mit Konfidenzintervall berichten
  \ite signifikant ist nicht dasselbe wie groß, nicht signifikant nicht
    dasselbe wie null
\ei
\os{\newpage}

\item[{\rot\bf Power: Wie viele Personen brauche ich?}]~\\[-\lblskip]
\begin{center}
\begin{tabular}{@{}lccc@{}}
\toprule
\textbf{erwarteter Effekt} & $d=0{,}2$ & $d=0{,}5$ & $d=0{,}8$ \\
\midrule
$n$ \emph{pro Gruppe} & 394 & 64 & 26 \\
\bottomrule
\end{tabular}
\end{center}
\bi
  \ite $t$-Test, zwei Gruppen, $\alpha=0{,}05$ zweiseitig, Power $0{,}80$
  \ite mit 13 Personen pro Gruppe findet man einen mittleren Effekt nur
    in etwa einem von vier Versuchen
\ei
\nsl{Die Werte der Tabelle sind mit der nichtzentralen $t$-Verteilung
nachgerechnet; sie stimmen mit den Tafeln bei \textcite{cohen1988power}
überein. Die Power bei 13 Personen pro Gruppe und $d=0{,}5$ beträgt
nach derselben Rechnung etwa 0,23. Kleine Studien haben nicht nur ein
hohes Risiko, echte Effekte zu übersehen. Wenn sie doch einmal
\enquote{signifikant} werden, überschätzen sie die Effektgröße
typischerweise deutlich \parencite{button2013power}.}
\os{\newpage}

\item[{\rot\bf Ein Lehrstück aus dem Buch}]~\\[-\lblskip]
\bi
  \ite Feldstudie zum Becher-Monty-Hall-Video: 32 auswertbare Fälle in
    drei Armen (13 / 9 / 10)
  \ite am 16.~Juli zeigte das Staunen eine schöne Abstufung: 3,83 / 4,08
    / 4,28
  \ite am Ende lagen die Arme gleichauf: 3,89 / 3,76 / 3,87
  \ite stärkster Zusammenhang: \textbf{Vorwissen} -- wer das Problem
    kannte, lag zu 92\,\% richtig
\ei
\nsl{Die Zahlen stammen aus der Feldstudie im Buch (Kap.~12f). Der
anfängliche Staunen-Gradient war ein Artefakt kleiner Zwischenstände:
Hätte die Erhebung früher geendet, stünde ein schöner, falscher Befund in
der Welt. Das Buch zieht daraus die Konsequenz, Fallzahlen vorab zu
planen (\enquote{ohne mindestens 60 Personen pro Arm lohnt der Aufwand
nicht}) und nach Vorwissen geschichtet zu randomisieren. Die Studie
berichtet bewusst keine p-Werte, sondern Effektgrößen (Hedges' $g$) mit
Bootstrap-Intervallen. Für die Stunde ist das ein Beispiel, das nah an
den Studierenden liegt: Das Spiel, das sie gerade gespielt haben, war
Gegenstand einer echten Studie mit echten Schwächen.}
\os{\newpage}
\end{description}

%==============================================================================
\section{Bayes: Was beweist ein Befund?}
%==============================================================================
\nsl{Der Satz von Bayes beschreibt, wie sich die Glaubwürdigkeit einer
Hypothese $H$ durch eine Beobachtung $E$ verändert. Paul Horwich hat
daraus eine Theorie der Evidenz entwickelt, die das Buch ausführlich
darstellt (Kap.~5b): Eine Beobachtung bestätigt eine Hypothese, wenn sie
sie wahrscheinlicher macht, als sie vorher war
\parencite{horwich1982probability}. Folgt die Beobachtung aus der
Hypothese, so ist der Zuwachs umso größer, je unwahrscheinlicher die
Beobachtung ohne Kenntnis der Hypothese war. Riskante Vorhersagen, die
eintreffen, sind deshalb starke Evidenz; Vorhersagen, die ohnehin
eingetroffen wären, sind schwache.}
\begin{description}
\item[{\rot\bf Der Satz von Bayes}]~\\[-\lblskip]
\[
P(H\mid E) \;=\; \frac{P(E\mid H)\cdot P(H)}{P(E)}
\]
\bi
  \ite $P(H)$: Prior -- was ich vorher für plausibel hielt
  \ite $P(E\mid H)$: Likelihood -- wie gut erwartet $H$ die Daten?
  \ite $P(H\mid E)$: Posterior -- was ich nachher glauben sollte
\ei
\os{\newpage}

\item[{\rot\bf Der Beweiswert riskanter Vorhersagen}]~\\[-\lblskip]
\bi
  \ite folgt $E$ aus $H$, gilt
    $P(H\mid E)\,/\,P(H) = 1\,/\,P(E)$
  \ite \textbf{je unwahrscheinlicher die Vorhersage, desto stärker die
    Bestätigung}, wenn sie eintrifft
  \ite \enquote{Morgen geht die Sonne auf} trifft ein und beweist
    fast nichts
  \ite \enquote{Wechseln gewinnt doppelt so oft} war riskant: Die
    Strichliste hätte sie widerlegen können
\ei
\nsl{Horwich zeigt außerdem, dass nur Beobachtungen, die ein
Widerlegungsrisiko eingehen, vollen Beweiswert haben
\parencite{horwich1982probability}. Das schlägt die Brücke zu LE~5
(Falsifizierbarkeit) und zu LE~10: Ein Verfahren, das jedes Ergebnis
liefern kann, geht kein Risiko ein und beweist deshalb nichts.}
\os{\newpage}
\end{description}

%==============================================================================
\section{Warnung: Bedeutung aus Zufall}
%==============================================================================
\nsl{Die letzte Theoriephase ist eine kurze Demonstration, die zeigt, wie
leicht ein \enquote{auffälliger} Befund entsteht, wenn man nur oft genug
testet. Jede Person sagt zehn Münzwürfe ihres Nachbarn vorher. Wer acht
oder mehr Treffer hat, gilt als \enquote{auffällig}: Die Wahrscheinlichkeit
dafür beträgt bei reinem Raten $56/1024 \approx 5{,}5\,\%$, liegt also
etwa an der üblichen Signifikanzschwelle. Bei 30 Personen ist die
Wahrscheinlichkeit, dass mindestens eine Person \enquote{auffällig} wird,
$1-(1-0{,}0547)^{30} \approx 81\,\%$; im Mittel sind es 1,6 Personen.
Wird jemand \enquote{hellsichtig}, ist die Demonstration gelungen. Wird
niemand hellsichtig, gelingt sie auch, weil man dann ausrechnet, wie
selten das ist. Die ausführliche Fassung mit dem Vorhersage-Trick aus dem
Buch (Abschnitt~10d) folgt in LE~10. Dass solche Freiheiten in echten
Studien die Fehlerquote weit über die nominellen 5\,\% treiben, haben
\textcite{simmons2011falsepositive} gezeigt.}
\begin{description}
\item[{\rot\bf Wer von Ihnen ist hellsichtig?}]~\\[-\lblskip]
\bi
  \ite sagen Sie zehn Münzwürfe Ihres Nachbarn vorher, notieren Sie die
    Treffer
  \ite 8 oder mehr Treffer gelten als \enquote{auffällig} (bei Raten:
    5,5\,\%)
  \ite wer hat 8 oder mehr?
  \ite bei 30 Personen gibt es mit etwa 81\,\% mindestens eine
    \enquote{Hellseherin}
\ei
\os{\newpage}

\item[{\rot\bf Was ist hier passiert?}]~\\[-\lblskip]
\bi
  \ite viele Tests: Ein \enquote{Treffer} ist fast garantiert
  \ite wer nur den Treffer berichtet, macht aus Zufall einen Befund
  \ite das ist der Kern von \textbf{p-Hacking}
  \ite Gegenprobe: Die \enquote{Hellseherin} sagt noch einmal zehn Würfe
    vorher
  \ite mehr dazu in LE~10: der Vorhersage-Trick
\ei
\nsl{Die Gegenprobe ist die eigentliche Pointe und sollte nach
Möglichkeit live gemacht werden. Die \enquote{auffällige} Person rät beim
zweiten Durchgang mit hoher Wahrscheinlichkeit nur noch durchschnittlich.
Das ist Regression zur Mitte und zugleich die Logik der Replikation.}
\os{\newpage}
\end{description}

\begin{merkbox}
Ein kleiner p-Wert sagt, dass die Daten schlecht zur Nullhypothese
passen. Er sagt nicht, dass die eigene Hypothese stimmt, nicht, wie groß
der Effekt ist, und gar nichts mehr, wenn man ihn unter vielen Versuchen
ausgesucht hat.
\end{merkbox}
\os{\newpage}

%==============================================================================
\section{Übung in der Stunde}
%==============================================================================
\begin{description}
\item[{\rot\bf Datensatz-Analyse (Kleingruppen, 20 Min.)}]~\\[-\lblskip]
\bi
  \ite Datensatz auf dem Übungsblatt Teil~A: Pendelausschlag von 50 Personen
    (fiktiv)
  \ite welches Skalenniveau hat jede Variable?
  \ite Median, Quartile, Mittelwert, Spannweite bestimmen
  \ite Histogramm und Boxplot zeichnen: Gibt es Ausreißer?
  \ite unterscheiden sich Personen mit und ohne Vorwissen?
\ei
\os{\newpage}

\item[{\rot\bf Zeitungsstatistik lesen (Plenum, 10 Min.)}]~\\[-\lblskip]
\bi
  \ite Übungsblatt Teil~B, Aufgabe~3, oder eine aktuelle Meldung
  \ite Was wurde gemessen, an wem, an wie vielen?
  \ite Was wird behauptet: Beschreibung, Unterschied, Ursache?
  \ite Wird eine Effektgröße genannt oder nur \enquote{signifikant}?
  \ite Wie viele Vergleiche wurden vermutlich gerechnet?
\ei
\os{\newpage}
\end{description}

\begin{lehrhinweis}
Die vollständig nachgerechneten Kennwerte des Datensatzes stehen im
Erwartungshorizont am Ende von Teil~B des Übungsblatts. Drei Stellen sollte man in der
Besprechung gezielt aufgreifen:
(1)~Quartile hängen vom Berechnungsverfahren ab; Taschenrechner,
Tabellenkalkulation und Handrechnung liefern leicht verschiedene Werte.
Das ist kein Fehler, muss aber angegeben werden. (2)~Der Vergleich der
Vorwissensgruppen kippt je nach Kennwert: Der Mittelwert spricht für einen
Unterschied, der Median nicht. Ursache sind der Ausreißer und die größere
Streuung in der Gruppe ohne Vorwissen. (3)~Die Staunen-Werte sind
ordinal; ein Mittelwert wird oft trotzdem berechnet. Bei der
Zeitungsstatistik genügt es, wenn die Gruppe mindestens zwei der vier
Fragen auf der Folie beantworten kann und benennt, welche Information
fehlt.
\end{lehrhinweis}

%==============================================================================
\section{Lernziele und Aufgaben}
%==============================================================================
\begin{description}
\item[{\rot\bf Nach dieser Einheit können Sie}]~\\[-\lblskip]
\bi
  \ite das Monty-Hall-Problem mit bedingter Wahrscheinlichkeit erklären
  \ite deskriptive und schließende Statistik unterscheiden
  \ite Lage- und Streuungsmaße passend zum Skalenniveau berechnen und
    deuten
  \ite die Logik eines Signifikanztests erklären und typische
    Fehldeutungen des p-Werts benennen
  \ite Effektgröße, Power und Stichprobengröße in Beziehung setzen
\ei
\os{\newpage}

\item[{\rot\bf Bis zur nächsten Woche}]~\\[-\lblskip]
\bi
  \ite \textbf{Datensatz-Analyse} fertigstellen (Übungsblatt Teil~A,
    Aufgabe~2)
  \ite \textbf{Kritik}: eine Zeitungsmeldung finden, die eine Studie mit
    $p<0{,}05$ als \enquote{Beweis} darstellt; methodische
    Kurzeinschätzung in höchstens zehn Sätzen (Teil~B, Aufgabe~4)
  \ite \textbf{Tagebuch des Staunens}: Wann hat Ihr Gefühl für
    Wahrscheinlichkeit Sie zuletzt getäuscht?
  \ite für LE~9: Was an Ihrer Forschungsfrage lässt sich \emph{nicht}
    zählen?
\ei
\os{\newpage}
\end{description}

\handoutinput{handouts/LE08}
